\section{The odd-dimensional orthogonal endpoint over three elements}
\label{sec:orthogonal-three-odd}

\paragraph{Status and conventions.}
This is a separate new author derivation.  The complete written
proof and exact supplement pass the author check; complete
proof reviews by /root and /root/g34\_eightfolds have also passed,
with a separate exact peer supplement.  It does not
change D-OL-12--13.  Let $S_d=\Omega_d(3)$ for odd $d\ge1135$,
and consider ordinary nontrivial irreducible complex modules
of this natural group.  Set
\[
 \beta:=\frac16-\frac1{10000}=\frac{4997}{30000},\qquad
 \delta_*:=\frac{1-(19/20)^9}{2}.
                                                               \tag{O3.1}
\]
A negative reflection means the actual natural isometry
$-r_v$, where $v$ is nonsingular and
$r_v(x)=x-B(x,v)Q(v)^{-1}v$.  Only those negative reflections
which belong to $S_d$ will be used.

\begin{lemma}[D-OL-14: the only low-support inner operators]
\label{D-OL-14}
Let $\rho$ be a nontrivial ordinary irreducible $S_d$-module of
degree $a$.  Every nonidentity $g\in S_d$ satisfies
\[
                          \nu(\rho(g))\ge\beta a.           \tag{O3.2}
\]
Unless $g$ is an actual negative reflection, the stronger bound
\[
                          \nu(\rho(g))>\delta_*a            \tag{O3.3}
\]
holds.  If $\rho$ is nonselfdual, every compatible linear
normalizer is a scalar times $\rho(g)$ for some $g\in S_d$.
Thus a nonscalar compatible normalizer with support at most
$\delta_*a$ is a scalar times a negative-reflection operator.
\end{lemma}

\begin{proof}
First take $g$ of prime order $l$.  If it fixes a hyperbolic
pair, use the opposite line radicals of D-OL-06.  Their fixed
space has dimension at most $a/2$.  On the complementary
natural space $W$, of dimension $N=d-2\ge20$, the moving rank
is at least two: a rank-one isometry in odd characteristic is
a reflection and has determinant $-1$, whereas $g\in S_d$.
The arbitrary-rank Gauss calculation (ON.3), as extended in
D-OL-10, bounds the fixed fraction in every nonzero norm orbit
by
\[
 \frac{3^{-2}+3^{-9}}{1-3^{-9}-3^{-19}}<\frac18.
                                                               \tag{O3.4}
\]
Its proof applies for every moving rank: up to half the dimension
the polar-rank bound gives Gauss error at most $3^{N/2}$; above
half the dimension the trivial error is already no larger.
Prime-order weight orbits therefore give
\[
               \nu(\rho(g))\ge\frac12\frac78\frac a2
                              =\frac7{32}a>\delta_*a.       \tag{O3.5}
\]

If $l$ is odd and no hyperbolic pair is fixed, the restricted
natural $\mu_2$-support is at least $567$.  Indeed, if $l\ne3$
the fixed space is nondegenerate of dimension at most two and
there is no eigenvalue $-1$, giving support at least $d-2$.
If $l=3$, put $r=\operatorname{rank}(g-I)$.  The fixed-space
polar rank is at least $d-2r$; absence of a hyperbolic pair
forces it to be at most two, so $r\ge(d-2)/2$.  As $d$ is
odd and at least $1135$, the integer rank is at least $567$.
Again $-1$ is not an eigenvalue.  Applying the exact
Larsen--Shalev--Tiep theorem and its restricted-support remark,
with the same ordinary Spin-to-$\Omega$ inflation and threshold
$567=7\cdot9^2$ as in (OO.8)--(OO.10), gives (O3.3).
The primary source is
\cite[Theorem~4.3.6, (4.3.2), pp.~1917--1918; complete proof
pp.~1918--1923; Remark~4.3.7, p.~1923]{HOS:LST2011}.

It remains to consider a prime-order involution with no fixed
hyperbolic pair.  Its $1$-eigenspace is nondegenerate of
dimension at most two.  Since $d$ is odd and the determinant
is one, that dimension is odd.  It is therefore one.  Its
orthogonal complement is the $-1$-eigenspace, so the involution
is precisely $-r_v$ for a nonsingular vector spanning its
$1$-eigenspace.

Such a negative reflection has a hyperbolic pair in its
$-1$-eigenspace.  In the signed coordinates (OO.3)--(OO.5),
the action on line-radical weight vectors is
$A=-g|_W$, an actual rank-one reflection.  The rank-one Gauss
bound, without replacing it by the false quarter bound, gives
\[
 \frac{|\mathcal O_t\cap\ker(A-I)|}{|\mathcal O_t|}
 \le c_3:=\frac{3^{-1}+3^{-9}}{1-3^{-9}-3^{-19}}
                  <1-4\beta.
                                                               \tag{O3.6}
\]
Together with the half-dimensional radical fixed-space bound
and the two-element weight orbits, this yields
\[
              \nu(\rho(g))\ge\frac{1-c_3}{4}a>\beta a.
                                                               \tag{O3.7}
\]
Thus all prime-order cases have been treated, with (O3.3)
failing to be asserted only for actual negative reflections.

We now classify arbitrary orders, rather than merely asserting
that some prime-order power is exceptional.  If the order of
$g$ has an odd prime divisor, a power of odd prime order is
already covered by (O3.3), and taking a power cannot increase
projective support.  Otherwise $g$ has order $2^s$.  If $s=1$,
the preceding involution analysis applies.  If $s\ge2$, put
$h=g^{2^{s-1}}$.  If $h$ is not a negative reflection, it
already gives the strong bound for $g$.

In the remaining case $h=-r_v$.  Since $g$ commutes with $h$,
it preserves the positive line $Fv$ and its perpendicular
$W_0$, on which $h=-I$.  A vector in a $1$- or $-1$-eigenspace
of $g$ is fixed by the even power $g^{2^{s-1}}$, and hence
belongs to $Fv$.  Thus both natural $\mu_2$-eigenspaces of
$g$ have dimension at most one.  The restricted support of
$g$ itself is at least $d-1\ge567$.  The published character
estimate applies directly to this element, without a
prime-order hypothesis, and proves (O3.3).  This proves the
claimed classification for every order.

Finally suppose $\rho$ is nonselfdual.  The normalizer argument
of (OO.11) uses only odd characteristic, not $q\ge5$.
For clarity, at $q=3$ there is no field automorphism; all
automorphisms in this large odd natural dimension are induced
by the full natural linear similitude group, by D-OL-07 and
\cite[Section~2.2.2, Table~2.1, p.~14]{HES:LiXia2022}.
In odd dimension that group is $F^*O(V)$ and
$O(V)=\langle-I\rangle\times SO(V)$, with
$SO(V)/S_d\simeq C_2$.

If the nontrivial linear outer coset stabilized $\rho$, all
of $O(V)$ would do so.  An irreducible constituent of
$\operatorname{Ind}_{S_d}^{O(V)}\rho$ would restrict to a
positive multiple of $\rho$.  Every irreducible character of
$O(V)$ is real-valued, by
\cite[Theorem~1, p.~44; exact domain and attribution,
pp.~39--40]{OL:KTV2023}.  This would make $\rho$ selfdual.
Hence the compatible linear normalizers are exactly scalar
times inner operators, as asserted.  This argument makes no
claim that all characters of $\Omega_d(3)$ are real.
\end{proof}

\begin{theorem}[D-OL-15: the remaining large odd-dimensional endpoint]
\label{D-OL-15}
A finite inducing image with an irreducibly acting nonempty normal central
product of nonselfdual ordinary $\Omega_{d_i}(3)$-modules,
all $d_i$ odd and at least $1135$, cannot supply period-image
rank nine for its actual six-dimensional multiplicity variation.
Mixed natural dimensions, scalars and all compatible factor
permutations are allowed; transitivity is not needed for this
assertion.

Consequently, in the main imaginary-quadratic sixfold problem,
every transitive ordinary orthogonal normal core of odd natural
dimension over any odd field satisfies
\[
                              d\le1133.                    \tag{O3.8}
\]
For selfdual factor modules the earlier bound $d\le21$ remains
valid.  These are necessary parameter bounds, not existence
claims at the remaining parameters.
\end{theorem}

\begin{proof}
Write $e=\prod_i a_i$ for the full inducing degree.  The
ordinary degree estimate of D-OL-01, or directly
\cite[Theorem~3.7, pp.~165--167]{OL:GLT2024}, gives
$a_i\ge2\cdot3^{d_i-4}>360$.
By D-OL-14 and tensor convolution, every genuine inertia has
actual moving rank at least $\beta e$.  Nontrivial factor
permutations have the stronger projective support $e/4$;
a nonidentity scalar has moving rank $e$.

Call an inertia good if its actual moving rank is at least
$\delta_*e$, and bad otherwise.  We record precisely what a
bad operator can be.  It cannot permute factors, since
$\delta_*<1/4$.  Write a factor-preserving operator as a
scalar times $\bigotimes_i\rho_i(g_i)$, using the complete
compatible-normalizer assertion in D-OL-14.  If any active
$g_i$ is not a negative reflection, tensor convolution gives
support greater than $\delta_*e$.  Hence every active factor
of a bad operator is a negative reflection.

Each such involution has two eigenspaces, with proportions
$p_i$ and $1-p_i$ lying in $[\beta,1-\beta]$.  For two active
factors, each of the two tensor eigenvalues has proportion at
most
\[
  \max_{p,q\in[\beta,1-\beta]}
       \{pq+(1-p)(1-q),\ p(1-q)+(1-p)q\}
       =(1-\beta)^2+\beta^2.
                                                               \tag{O3.9}
\]
The equality follows because each expression is affine in each
variable separately, so its maximum is at a corner of the
rectangle.  Therefore two active factors give support at least
$2\beta(1-\beta)e>e/4>\delta_*e$.  Additional convolution
cannot increase a largest eigenvalue proportion.  Every bad
inertia consequently has exactly one active factor and is a
scalar times one negative-reflection operator.  That scalar
is $1$ or $-1$: otherwise neither of its two eigenvalues is
one and the actual moving rank is $e$.  No general assertion
that scalar convolution preserves a nonreal-eigenvalue bound
is used here.

The exact arithmetic needed for the source budget is
\[
 13\beta-2>\frac17,\qquad
 11\beta+\delta_*-2>\frac1{60},\qquad
                  2\beta(1-\beta)>\frac14.                 \tag{O3.10}
\]
The actual intermediate-cover Euler formula is
$6=2e(h-1)+D$.  Since $e>360$, genus at least two is impossible.
Genus one has $D=6$, at most six genuine marks, and period-image
rank at most six.  Rank nine would require genus zero, at least
twelve genuine inertia points, and $D=2e+6$.

If their number is at least thirteen, then
$2e+6\ge13\beta e$, whence $e<42$ by (O3.10), a contradiction.
Thus there are exactly twelve.  If any one is good, then
$2e+6\ge(\delta_*+11\beta)e$, whence $e<360$, another
contradiction.  Therefore all twelve would be bad.

We now use actual generation, not only the preceding dimension
and rank inequalities.  The twelve genus-zero inertia elements
generate the faithful inducing image $J$.  Since every bad
inertia preserves each factor, so does all of $J$.  Fix one
normal factor $S_{d_i}$.  The induced action of $J$ on that
factor is generated by the inner automorphisms of at most
twelve negative reflections $g_1,\ldots,g_k$, with $k\le12$;
inactive generators induce the identity.  But $J$ contains
the normal factor itself, so this induced image contains
all of $\operatorname{Inn}(S_{d_i})$.  As the natural group
$S_{d_i}$ has trivial center, the elements $g_j$ generate
$S_{d_i}$.

Write $g_j=-r_{v_j}$ in its natural quadratic space, and put
\[
              U=\bigcap_{j=1}^k v_j^\perp,
                     \qquad \dim U\ge d_i-k>0.             \tag{O3.11}
\]
Every $g_j$ restricts to $-I$ on $U$.  The group they generate
therefore acts on this nonzero space through the scalar
group $\{I,-I\}$, and at least one generator has nontrivial
image.  This is a nontrivial homomorphism from the perfect
group $S_{d_i}$ to a group of order two, a contradiction.
There is at least one active generator because the induced
image contains the nontrivial factor.  The natural dimension
bound in (O3.11) is much larger than needed, but is the same
one already used for the exact character estimate.

This excludes the nonselfdual source without using a
similitude character, a determinant class, or equality of the
inducing character field with the original imaginary field.
For the final main-source consequence, transitivity preserves
selfduality, so all factor modules are selfdual or all are
nonselfdual.  The former case is excluded at $d\ge22$ by
D-OL-09.  In the latter case $q=3$ is excluded at $d\ge1135$
by this theorem, and every odd $q\ge5$ is excluded at that
threshold by D-OL-12.  Since $d$ is odd, (O3.8) follows.
\end{proof}

\paragraph{Certificate and scope.}
The exact supplement checks the rational constants, the
two-involution convolution corners, an actual higher-order
root of a negative reflection, and a twelve-reflection
common-perpendicular model.  The all-order classification,
compatible-normalizer reduction and generation contradiction
are the complete arguments above, not extrapolations from
these examples.  Ordinary modules are essential: no claim
about arbitrary noninflated spin characters is added here.
